
\section{Conclusion}
\label{sec:conclusion}

This paper treated the choice among established urban routing policies as a
per-instance decision problem and measured it on a controlled factorial of 648
operating contexts, four policies and five paired seeds, comprising $12{,}960$
valid microsimulation runs with no failures and no teleports.

The portfolio is complementary. Three of four policies are strictly best in at
least one context with a margin that survives replication, the reactive
travel-time policy is never a resolved winner, and 73.0\,\% of contexts have more
than one non-dominated policy across journey time, \CO\ and stopped delay. The
winner map is genuinely multi-factor: no single factor predicts it above
66.5\,\%, and the organising quantity is the product of demand and guidance
penetration relative to the capacity of the shortest corridor.

The decision value of that complementarity is small. The single best fixed policy
leaves 0.45\,s per vehicle of headroom on a mean journey time of 474.4\,s, because
it attains the per-context minimum in two thirds of contexts and, where it does
not, loses 1.32\,s on average against winning by 22.47\,s where it does. By
contrast, committing to the worst fixed policy costs 220.2\,s per vehicle, 46.4\,\%
of the best fixed policy's mean. The first-order choice is worth roughly 490 times
the second-order one, and tuning a single policy parameter was worth 37 times the
entire adaptation headroom.

Deployability is a third and separate question. Under leave-one-context-out
evaluation a gradient-boosted model trained on advantage recovers 55.8\,\% of the
headroom and more than halves the worst case, while a transparent depth-2 rule is
more accurate than doing nothing and costs 71\,\% more; accuracy and decision
quality order the candidate rules differently. Across eight distribution shifts
reported separately by type, the selector helps on two, is indistinguishable on
three and harms on three, with the worst damage under concept shift. A
support-and-confidence gate never certified a deviation in any of the 648
contexts, because the conformal interval is two orders of magnitude wider than the
effect, and its support component flagged every context of three factor-level
extrapolations but none of the domain shift, which its feature representation
cannot see.

The contribution is therefore a measurement and a set of distinctions, not an
algorithm. Complementarity, decision value and deployability are different
properties; they are quantitatively disconnected in this environment; and a study
that reports the first as evidence for the third --- as is common when portfolios
are described by how many members win somewhere --- would overstate the case for
adaptive routing-policy selection by more than two orders of magnitude. The
practical recommendation that follows is correspondingly modest: choose the fixed
policy carefully, calibrate it, and add an adaptive layer only after measuring
that the headroom it competes for exceeds the noise of the system it will act on.

All results are scoped to one synthetic network and are not externally validated.
The specification, run records, analysis code and figure scripts are available so
that every number can be recomputed from the simulation output.
